% Compilar dos veces con pdflatex desde este directorio.
\documentclass[aspectratio=169,10pt]{beamer}
\usepackage[utf8]{inputenc}
\usepackage[T1]{fontenc}
\usepackage[spanish,es-nodecimaldot]{babel}
\usepackage{amsmath,amssymb,booktabs,tabularx,array,tikz}
\usetikzlibrary{arrows.meta,positioning,calc,shapes.geometric}

\definecolor{UAPNavy}{HTML}{17324D}
\definecolor{UAPTeal}{HTML}{007F82}
\definecolor{UAPGold}{HTML}{E2A33A}
\definecolor{UAPLight}{HTML}{F3F6F8}
\definecolor{UAPGray}{HTML}{536471}
\definecolor{UAPRed}{HTML}{A94343}

\tikzset{
  flow/.style={-{Stealth[length=2.2mm]},thick,draw=UAPNavy},
  datum/.style={draw=UAPGold,fill=UAPGold!16,rounded corners=2pt,align=center,minimum height=8mm,text=UAPNavy},
  process/.style={draw=UAPTeal,fill=UAPTeal!12,rounded corners=3pt,very thick,align=center,minimum height=10mm,text=UAPNavy},
  decision/.style={draw=UAPNavy,fill=UAPNavy!9,rounded corners=3pt,very thick,align=center,minimum height=10mm,text=UAPNavy},
  warning/.style={draw=UAPRed,fill=UAPRed!7,rounded corners=3pt,align=left}
}

\usetheme{Boadilla}
\usecolortheme{default}
\setbeamercolor{normal text}{fg=UAPNavy,bg=white}
\setbeamercolor{structure}{fg=UAPTeal}
\setbeamercolor{frametitle}{fg=white,bg=UAPNavy}
\setbeamercolor{title}{fg=white}
\setbeamercolor{subtitle}{fg=UAPGold}
\setbeamercolor{block title}{fg=white,bg=UAPTeal}
\setbeamercolor{block body}{fg=UAPNavy,bg=UAPLight}
\setbeamercolor{alertblock title}{fg=white,bg=UAPRed}
\setbeamercolor{alertblock body}{fg=UAPNavy,bg=UAPRed!7}
\setbeamerfont{frametitle}{series=\bfseries}
\setbeamertemplate{navigation symbols}{}
\setbeamertemplate{itemize item}{\color{UAPGold}$\blacktriangleright$}
\setbeamertemplate{itemize subitem}{\color{UAPTeal}$\bullet$}
\setbeamertemplate{footline}{%
  \leavevmode\hbox{%
    \begin{beamercolorbox}[wd=.78\paperwidth,ht=2.8ex,dp=1.2ex,leftskip=1em]{author in head/foot}%
      \color{white}\insertshorttitle
    \end{beamercolorbox}%
    \begin{beamercolorbox}[wd=.22\paperwidth,ht=2.8ex,dp=1.2ex,rightskip=1em plus1fil]{date in head/foot}%
      \color{UAPGray}\hfill Semana 8 \enspace | \enspace \insertframenumber/\inserttotalframenumber
    \end{beamercolorbox}}}

\newcommand{\concepto}[1]{\textcolor{UAPTeal}{\textbf{#1}}}
\newcommand{\pregunta}[1]{\begin{block}{Pregunta clave}#1\end{block}}
\title[IA · Clase 8]{Aprendizaje no supervisado}
\subtitle{De la semejanza a segmentos e hipótesis}
\author{Curso de Inteligencia Artificial}
\institute{Semana 8 · Clase 8}
\date{}

\begin{document}

% 1
\begin{frame}[plain]
  \begin{beamercolorbox}[wd=\paperwidth,ht=\paperheight,sep=2em]{frametitle}
    \vfill
    {\Huge\bfseries\inserttitle\par}\vspace{.45em}
    {\color{UAPGold}\Large\insertsubtitle\par}\vspace{1.2em}
    \begin{tikzpicture}[node distance=5mm,scale=.78,transform shape]
      \node[datum] (r) {Representación};
      \node[process,right=of r] (s) {Semejanza};
      \node[datum,right=of s] (e) {Estructura};
      \node[decision,right=of e] (i) {Interpretación};
      \draw[flow] (r)--(s); \draw[flow] (s)--(e); \draw[flow] (e)--(i);
    \end{tikzpicture}\par
    \vfill {\color{UAPGold}\large\insertauthor\par}{\color{white}\insertinstitute\par}
  \end{beamercolorbox}
\end{frame}

% 2
\begin{frame}{Propósito de la clase}
  Comprender cómo los métodos no supervisados organizan observaciones sin una etiqueta objetivo y cómo evaluar si esa organización es útil.
  \begin{block}{Al finalizar, podremos}
    \begin{itemize}
      \item formular una pregunta no supervisada;
      \item justificar una distancia y una representación;
      \item explicar K-means, jerárquico, DBSCAN y SOM;
      \item usar PCA para resumir y visualizar datos;
      \item interpretar segmentos sin convertirlos en categorías esenciales.
    \end{itemize}
  \end{block}
\end{frame}

% 3
\begin{frame}{Punto de partida}
  \small
  En aprendizaje supervisado disponíamos de pares con respuesta observada:
  \begin{center}\concepto{$(x_i,y_i)$ $\rightarrow$ modelo ajustado $\rightarrow$ predicción}\end{center}
  Ahora solo observamos perfiles:
  \begin{center}\concepto{$x_1,\ldots,x_n$ $\rightarrow$ representación $\rightarrow$ estructura candidata}\end{center}
  \pregunta{¿Qué significa que dos observaciones sean semejantes para nuestro problema?}
\end{frame}

% 4
\begin{frame}{Caso conductor: movilidad}
  \small
  Una ciudad quiere resumir perfiles de zonas para planificar observación y servicio.
  \begin{tabularx}{\textwidth}{>{\bfseries}p{.27\textwidth}X}\toprule
  Elemento & Ejemplo\\\midrule
  Unidad & zona de movilidad durante una semana\\
  Formación & proporciones por franja, duración típica, variabilidad y congestión\\
  Sin etiqueta & no existe un tipo verdadero de zona\\
  Uso posterior & describir perfiles, revisar cobertura, proponer acciones exploratorias\\\bottomrule
  \end{tabularx}
  \begin{block}{Idea}Un grupo resume una estructura bajo decisiones explícitas; no descubre una categoría natural.\end{block}
\end{frame}

% 5
\begin{frame}{¿Qué cambia sin etiquetas?}
  \small
  \begin{tabularx}{\textwidth}{>{\bfseries}p{.26\textwidth}XX}\toprule
   & Supervisado & No supervisado\\\midrule
  Información & objetivo $y$ observado & sin respuesta correcta conocida\\
  Calidad & error frente a $y$ & cohesión, separación, estabilidad y utilidad\\
  Resultado & predicción para casos nuevos & segmentos, prototipos o vecindades\\
  Revisión & validación con etiquetas & evidencia interna, externa y contextual\\\bottomrule
  \end{tabularx}
  \begin{alertblock}{Consecuencia}Muchas particiones pueden ser válidas matemáticamente. Ninguna métrica única decide cuál sirve mejor.\end{alertblock}
\end{frame}

% 6
\begin{frame}{Cuatro objetivos posibles}
  \small
  \begin{itemize}
    \item \concepto{Descubrimiento de grupos:} semejanza interna y diferencia entre segmentos.
    \item \concepto{Representación:} resumir con prototipos, componentes o mapas.
    \item \concepto{Estructura local:} identificar vecindades, regiones densas y puntos aislados.
    \item \concepto{Generación de hipótesis:} proponer preguntas para contrastar después.
  \end{itemize}
  \begin{alertblock}{Control}Elegir un algoritmo sin declarar el objetivo confunde una figura atractiva con un resultado útil.\end{alertblock}
\end{frame}

% 7
\begin{frame}{Formular antes de agrupar}
  \small
  Antes del ajuste debemos responder:
  \begin{enumerate}
    \item ¿Cuál es la unidad de análisis?
    \item ¿Qué periodo y población cubren los datos?
    \item ¿Qué diferencias importan para el dominio?
    \item ¿Se admite ruido sin asignar a un grupo?
    \item ¿Cómo se usará y revisará el resultado?
  \end{enumerate}
  \begin{block}{Idea}La distancia es una hipótesis de dominio expresada matemáticamente.\end{block}
\end{frame}

% 8
\begin{frame}{Distancia euclídea}
  \small
  Para perfiles numéricos $\mathbf{x}$ y $\mathbf{z}$:
  \[d_2(\mathbf{x},\mathbf{z})=\sqrt{\sum_{j=1}^{p}(x_j-z_j)^2}.\]
  Para $A=(2,3)$ y $B=(5,7)$:
  \[d_2(A,B)=\sqrt{3^2+4^2}=5.\]
  \begin{itemize}
    \item Las diferencias grandes pesan más al elevarse al cuadrado.
    \item Es coherente con K-means y centroides definidos por medias.
    \item Supone variables cuantitativas y escalas comparables.
  \end{itemize}
\end{frame}

% 9
\begin{frame}{Distancia Manhattan}
  \small
  La distancia Manhattan suma diferencias absolutas:
  \[d_1(\mathbf{x},\mathbf{z})=\sum_{j=1}^{p}|x_j-z_j|.\]
  Para los mismos perfiles:
  \[d_1(A,B)=|2-5|+|3-7|=7.\]
  \begin{block}{Lectura}No elimina la necesidad de escalar variables, pero reduce el peso relativo de diferencias extremas frente a $d_2$.\end{block}
\end{frame}

% 10
\begin{frame}{La métrica puede cambiar el vecino}
  \small
  Para $O=(0,0)$, $P=(3,0)$ y $Q=(2,2)$:
  \begin{tabularx}{\textwidth}{>{\bfseries}Xccc}\toprule
  Métrica & $d(O,P)$ & $d(O,Q)$ & Vecino más cercano\\\midrule
  Euclídea & $3$ & $\sqrt{8}\approx2{,}83$ & $Q$\\
  Manhattan & $3$ & $4$ & $P$\\\bottomrule
  \end{tabularx}
  \begin{alertblock}{Conclusión}No hay una métrica universalmente correcta: cada una expresa una noción distinta de diferencia.\end{alertblock}
\end{frame}

% 11
\begin{frame}{Similitud coseno: forma y nivel}
  \small
  La similitud coseno compara orientación:
  \[s_{\cos}(\mathbf{x},\mathbf{z})=\frac{\mathbf{x}^{\mathsf T}\mathbf{z}}{\|\mathbf{x}\|_2\|\mathbf{z}\|_2}.\]
  Los perfiles $A=(2,4,2)$ y $B=(20,40,20)$ tienen coseno $1$: misma composición relativa, distinta magnitud.
  \begin{block}{Decisión}Usar proporciones para la forma y conservar el volumen como una variable adicional si importa operacionalmente.\end{block}
\end{frame}

% 12
\begin{frame}{Escalar también es decidir}
  \small
  Una duración entre 0 y 120 puede dominar a una proporción entre 0 y 1. Dos opciones frecuentes:
  \[z_{ij}=\frac{x_{ij}-\bar x_j}{s_j},\qquad r_{ij}=\frac{x_{ij}-\operatorname{mediana}(X_j)}{Q_{0.75}(X_j)-Q_{0.25}(X_j)}.\]
  \begin{itemize}
    \item $z$: centra y asigna varianza unitaria.
    \item $r$: alternativa robusta basada en mediana y rango intercuartílico.
    \item Los logaritmos pueden comprimir colas derechas.
  \end{itemize}
\end{frame}

% 13
\begin{frame}{Ponderar bloques de variables}
  \small
  Una distancia ponderada explicita prioridades:
  \[d_w(\mathbf{x},\mathbf{z})=\sqrt{\sum_{j=1}^{p}w_j(x_j-z_j)^2},\qquad w_j\geq0.\]
  En movilidad, doce proporciones horarias pueden dominar por cantidad de columnas a dos variables de diversidad, incluso estandarizadas.
  \begin{block}{Criterio}Ponderar por confiabilidad, relevancia operativa o equilibrio entre bloques; no para producir grupos más vistosos.\end{block}
\end{frame}

% 14
\begin{frame}{Datos mixtos y variables reservadas}
  \small
  \begin{itemize}
    \item K-means requiere variables numéricas y una media interpretable.
    \item Las categorías necesitan una codificación y una disimilitud coherentes.
    \item Las variables de \concepto{formación} determinan los grupos.
    \item Las variables de \concepto{caracterización} se reservan para comprenderlos después.
  \end{itemize}
  \begin{alertblock}{Control}Separar ambos conjuntos evita introducir previamente la conclusión que se busca describir.\end{alertblock}
\end{frame}

% 15
\begin{frame}{K-means: idea y objetivo}
  \small
  K-means divide observaciones numéricas en $K$ grupos $C_1,\ldots,C_K$ minimizando:
  \[J=\sum_{k=1}^{K}\sum_{\mathbf{x}_i\in C_k}\|\mathbf{x}_i-\boldsymbol{\mu}_k\|_2^2,\qquad \boldsymbol{\mu}_k=\frac{1}{|C_k|}\sum_{\mathbf{x}_i\in C_k}\mathbf{x}_i.\]
  \begin{block}{Lectura}El centroide es un perfil promedio; no tiene por qué coincidir con una zona real.\end{block}
\end{frame}

% 16
\begin{frame}{K-means: dos pasos alternados}
  \small
  \begin{enumerate}
    \item Inicializar $K$ centroides.
    \item Asignar cada observación al centroide más cercano.
    \item Recalcular cada centroide como media de su grupo.
    \item Repetir hasta que no cambien asignaciones o la mejora sea pequeña.
  \end{enumerate}
  \begin{itemize}
    \item Cada paso no aumenta $J$.
    \item Converge a una solución local, no garantiza el óptimo global.
    \item Un grupo vacío requiere una regla documentada.
  \end{itemize}
\end{frame}

% 17
\begin{frame}{Qué estructura busca K-means}
  \small
  Funciona mejor con grupos:
  \begin{columns}[T]
    \column{.48\textwidth}
    \begin{block}{Adecuados}\begin{itemize}\item compactos;\item aproximadamente esféricos;\item de dispersión y tamaño semejantes.\end{itemize}\end{block}
    \column{.48\textwidth}
    \begin{alertblock}{Problemáticos}\begin{itemize}\item alargados;\item densidades muy distintas;\item observaciones que deberían quedar sin grupo.\end{itemize}\end{alertblock}
  \end{columns}
  \begin{block}{Límite}Su eficiencia no corrige una geometría inadecuada.\end{block}
\end{frame}

% 18
\begin{frame}{Inicialización, reinicios y estabilidad}
  \small
  \begin{itemize}
    \item Distintas semillas pueden producir soluciones diferentes.
    \item K-means++ favorece centros iniciales separados.
    \item Varios reinicios reducen el riesgo de una mala solución local.
    \item Comparar solo la menor inercia puede ocultar asignaciones muy distintas.
  \end{itemize}
  \begin{block}{Control}Además de cambiar semillas, repetir con submuestras, periodos y transformaciones plausibles.\end{block}
\end{frame}

% 19
\begin{frame}{Elegir el número de grupos}
  \small
  La inercia disminuye siempre al aumentar $K$. El codo busca rendimientos decrecientes, pero puede no ser claro.
  \begin{enumerate}
    \item Definir un rango de $K$ útil para el problema.
    \item Ejecutar varios reinicios por cada $K$.
    \item Comparar inercia, silhouette, Davies--Bouldin y tamaños.
    \item Revisar estabilidad y perfiles.
    \item Elegir la solución más simple que sea estable y accionable.
  \end{enumerate}
\end{frame}

% 20
\begin{frame}{Clustering jerárquico}
  \small
  Construye relaciones anidadas entre observaciones.
  \begin{tabularx}{\textwidth}{>{\bfseries}p{.22\textwidth}XX}\toprule
  Estrategia & Inicio & Operación\\\midrule
  Aglomerativa & una observación por grupo & fusiona los grupos más próximos\\
  Divisiva & todas las observaciones juntas & divide grupos sucesivamente\\\bottomrule
  \end{tabularx}
  \begin{block}{Resultado}El dendrograma permite examinar varios niveles de granularidad sin fijar $K$ al inicio.\end{block}
\end{frame}

% 21
\begin{frame}{Reglas de enlace}
  \small
  La distancia entre grupos $A$ y $B$ depende de la regla elegida:
  \[D_{\text{simple}}(A,B)=\min_{i\in A,\,l\in B}d(i,l),\qquad D_{\text{completo}}(A,B)=\max_{i\in A,\,l\in B}d(i,l).\]
  \[D_{\text{promedio}}(A,B)=\frac{1}{|A||B|}\sum_{i\in A}\sum_{l\in B}d(i,l).\]
  \begin{itemize}
    \item Simple: descubre cadenas, pero puede unir regiones por puntos puente.
    \item Completo: favorece grupos compactos; es sensible a extremos.
    \item Promedio: usa toda la información entre pares.
    \item Ward: minimiza aumento de dispersión y requiere geometría euclídea coherente.
  \end{itemize}
\end{frame}

% 22
\begin{frame}{Leer un dendrograma}
  \small
  \begin{itemize}
    \item Cada hoja es una observación o grupo inicial.
    \item La altura de unión indica la disimilitud de enlace.
    \item Un corte horizontal induce una partición.
  \end{itemize}
  \begin{alertblock}{Tres errores frecuentes}
  Leer la posición horizontal como distancia; elegir un corte solo por cantidad deseada; tratar ramas como categorías reales sin perfiles ni estabilidad.
  \end{alertblock}
  \begin{block}{Idea}El dendrograma explora granularidades; no cuenta grupos automáticamente.\end{block}
\end{frame}

% 23
\begin{frame}{Densidad y vecindades}
  \small
  Una vecindad de radio $\varepsilon$ alrededor de $\mathbf{x}$ es:
  \[N_{\varepsilon}(\mathbf{x})=\{\mathbf{z}:d(\mathbf{x},\mathbf{z})\leq\varepsilon\}.\]
  El tamaño $|N_{\varepsilon}(\mathbf{x})|$ aproxima densidad local.
  \begin{itemize}
    \item Los métodos por densidad encuentran regiones conectadas y densas.
    \item Pueden describir formas no convexas.
    \item Pueden dejar casos sin grupo.
  \end{itemize}
\end{frame}

% 24
\begin{frame}{DBSCAN: núcleo, frontera y ruido}
  \small
  DBSCAN usa radio $\varepsilon$ y \texttt{minPts}:
  \begin{tabularx}{\textwidth}{>{\bfseries}p{.22\textwidth}X}\toprule
  Tipo & Condición\\\midrule
  Núcleo & su vecindad contiene al menos \texttt{minPts} casos\\
  Frontera & no es núcleo, pero está cerca de uno\\
  Ruido & no cumple ninguna de las condiciones anteriores\\\bottomrule
  \end{tabularx}
  \begin{block}{Mecanismo}Expande grupos mediante cadenas de puntos núcleo conectados por densidad.\end{block}
\end{frame}

% 25
\begin{frame}{DBSCAN frente a K-means}
  \small
  \begin{tabularx}{\textwidth}{>{\bfseries}p{.25\textwidth}XX}\toprule
   & K-means & DBSCAN\\\midrule
  Parámetros & $K$ & $\varepsilon$ y \texttt{minPts}\\
  Asignación & todos los casos & puede dejar ruido\\
  Geometría & grupos convexos & formas arbitrarias\\
  Resumen & centroide promedio & regiones densas conectadas\\\bottomrule
  \end{tabularx}
  \begin{alertblock}{Precaución}Ruido es aislamiento geométrico bajo una configuración; no equivale a error, fraude ni evento irrelevante.\end{alertblock}
\end{frame}

% 26
\begin{frame}{Sensibilidad de DBSCAN}
  \small
  \begin{itemize}
    \item \texttt{minPts} pequeño acepta cadenas frágiles; uno grande convierte más casos en ruido.
    \item La distancia al vecino $k$ puede sugerir un rango para $\varepsilon$, no decidirlo por sí sola.
    \item El escalado cambia el significado efectivo de $\varepsilon$.
  \end{itemize}
  \begin{block}{Control}Explorar una rejilla de parámetros y buscar mesetas donde cantidad, tamaño y composición de grupos cambien poco.\end{block}
\end{frame}

% 27
\begin{frame}{Mapas autoorganizados (SOM)}
  \small
  Un SOM es una rejilla de prototipos. Cada nodo $m$ tiene pesos $\mathbf{w}_m$.
  Para una observación $\mathbf{x}$, la unidad de mejor coincidencia es:
  \[c=\arg\min_m\|\mathbf{x}-\mathbf{w}_m\|_2.\]
  Durante el ajuste se acercan a $\mathbf{x}$ la unidad ganadora y sus vecinas de rejilla:
  \[\mathbf{w}_m(t+1)=\mathbf{w}_m(t)+\alpha(t)h_{cm}(t)[\mathbf{x}(t)-\mathbf{w}_m(t)].\]
\end{frame}

% 28
\begin{frame}{Qué muestra un SOM}
  \small
  \begin{itemize}
    \item \concepto{Mapa de componentes:} valor de una variable en cada prototipo.
    \item \concepto{Matriz U:} distancia entre prototipos vecinos; sugiere fronteras o continuidad.
    \item \concepto{Mapa de conteos:} cantidad de observaciones por nodo.
  \end{itemize}
  \begin{block}{Límite}Un SOM preserva vecindades de forma aproximada; no produce grupos automáticamente. Agrupar prototipos es una etapa adicional.\end{block}
\end{frame}

% 29
\begin{frame}{Reducir dimensión: por qué}
  \small
  Con muchas variables, las distancias tienden a parecerse y las vecindades pierden contraste.
  \[Z=g(X),\qquad q<p.\]
  \begin{itemize}
    \item Puede comprimir, disminuir ruido o facilitar visualización.
    \item Puede ocultar señales raras o distorsionar proximidades.
    \item No debe justificarse solo porque produce una figura más clara.
  \end{itemize}
  \begin{alertblock}{Control}No agrupar solo una proyección 2D sin revisar la estructura en el espacio de formación.\end{alertblock}
\end{frame}

% 30
\begin{frame}{PCA: componentes principales}
  \small
  PCA construye nuevas coordenadas lineales que resumen variación de $X$.
  \begin{enumerate}
    \item Centrar las variables.
    \item Decidir si estandarizar según las unidades.
    \item Buscar direcciones ortogonales de máxima varianza.
    \item Proyectar cada observación en esas direcciones.
  \end{enumerate}
  \begin{block}{Lectura}La primera componente explica la mayor varianza disponible; las siguientes explican varianza restante y son ortogonales a las anteriores.\end{block}
\end{frame}

% 31
\begin{frame}{Varianza explicada e interpretación}
  \small
  Si $\lambda_j$ es la varianza de la componente $j$:
  \[\text{varianza explicada}_j=\frac{\lambda_j}{\sum_r\lambda_r}.\]
  \begin{itemize}
    \item Una componente combina variables originales mediante cargas.
    \item Alta varianza no implica utilidad para todas las decisiones.
    \item PCA optimiza reconstrucción lineal, no clasificación, causalidad ni calidad de grupos.
  \end{itemize}
  \begin{block}{Resultado}Elegir $q$ combina pérdida aceptable, estabilidad e interpretación.\end{block}
\end{frame}

% 32
\begin{frame}{Inercia: compactación de K-means}
  \small
  La inercia es la función objetivo:
  \[W_K=\sum_{k=1}^{K}\sum_{i\in C_k}\|\mathbf{x}_i-\boldsymbol{\mu}_k\|_2^2.\]
  \begin{itemize}
    \item Menor valor: menor dispersión bajo esa representación.
    \item Solo compara soluciones sobre los mismos datos y preprocesamiento.
    \item Siempre disminuye al crecer $K$; no basta para elegirlo.
  \end{itemize}
\end{frame}

% 33
\begin{frame}{Silhouette: cohesión y separación}
  \small
  Para $i$, $a(i)$ es distancia media a su grupo y $b(i)$ la menor distancia media a otro grupo:
  \[s(i)=\frac{b(i)-a(i)}{\max\{a(i),b(i)\}}.\]
  \begin{tabularx}{\textwidth}{>{\bfseries}p{.28\textwidth}X}\toprule
  Valor & Lectura\\\midrule
  cercano a 1 & buena cohesión y separación\\
  cercano a 0 & punto de frontera\\
  negativo & más cercano, en promedio, a otro grupo\\\bottomrule
  \end{tabularx}
  \begin{block}{Control}Revisar la distribución por grupo, no solo su promedio.\end{block}
\end{frame}

% 34
\begin{frame}{Davies--Bouldin: grupos rivales}
  \small
  Compara dispersión interna y distancia entre centroides:
  \[DB=\frac{1}{K}\sum_{k=1}^{K}\max_{l\neq k}\frac{S_k+S_l}{M_{kl}}.\]
  Valores menores indican grupos compactos respecto de su rival más cercano.
  \begin{alertblock}{Límite}Supone grupos descritos por centroides; puede valorar mal formas arbitrarias y resultados de DBSCAN.\end{alertblock}
\end{frame}

% 35
\begin{frame}{Caracterizar sin esencializar}
  \small
  Una ficha de segmento debe indicar:
  \begin{itemize}
    \item tamaño, perfiles centrales y dispersión;
    \item variables que lo distinguen y casos fronterizos;
    \item estabilidad ante muestras, periodos y decisiones de representación;
    \item qué casos no representa;
    \item para qué pregunta o acción exploratoria podría servir.
  \end{itemize}
  \begin{block}{Lenguaje}Usar nombres descriptivos, como ``actividad nocturna alta y variabilidad baja'', en lugar de etiquetas valorativas.\end{block}
\end{frame}

% 36
\begin{frame}{Anomalías y sobreinterpretación}
  \small
  Una observación atípica puede ser global, contextual o colectiva. Primero se revisan integridad, cobertura, unidades y periodo.
  \begin{alertblock}{No debemos concluir que}
    \begin{itemize}
      \item un cluster es una clase natural;
      \item una separación visual prueba causalidad;
      \item ruido de DBSCAN es error;
      \item una métrica alta vuelve estable o útil la segmentación.
    \end{itemize}
  \end{alertblock}
  \begin{block}{Control}Triangular métricas internas, estabilidad, evidencia externa y utilidad contextual.\end{block}
\end{frame}

% 37
\begin{frame}{Del segmento a la decisión}
  \small
  En movilidad, los segmentos pueden guiar preguntas como:
  \begin{tabularx}{\textwidth}{>{\bfseries}p{.42\textwidth}X}\toprule
  Hallazgo descriptivo & Pregunta posterior\\\midrule
  patrón nocturno concentrado & ¿la cobertura de observación es suficiente?\\
  zona frontera entre segmentos & ¿cambia por periodo o por medición?\\
  perfil aislado & ¿es un evento raro, una falla o una zona subobservada?\\\bottomrule
  \end{tabularx}
  \begin{block}{Idea}El grupo informa una investigación o una regla revisable; no reemplaza el juicio ni demuestra por qué ocurre el patrón.\end{block}
\end{frame}

% 38
\begin{frame}{Síntesis}
  \small
  \begin{itemize}
    \item El aprendizaje no supervisado organiza perfiles sin una respuesta conocida.
    \item Representación, distancia y escalado definen qué estructura puede aparecer.
    \item K-means, jerárquico, DBSCAN y SOM responden a geometrías diferentes.
    \item PCA resume variación, pero no certifica utilidad ni conserva toda la estructura.
    \item La evaluación combina compactación, separación, estabilidad, evidencia externa y propósito.
    \item Un segmento es una descripción condicionada por método y datos, no una identidad fija.
  \end{itemize}
\end{frame}

% 39
\begin{frame}{Lecturas}
  \small
  \begin{block}{Lectura principal}
    \begin{itemize}
      \item Capítulo 13, §§13.1--13.4: distancias, clustering, DBSCAN, SOM y evaluación.
    \end{itemize}
  \end{block}
  \begin{block}{Para ampliar}
    \begin{itemize}
      \item Capítulo 10, §10.3: reducción de dimensionalidad, PCA y visualización.
    \end{itemize}
  \end{block}
  \begin{block}{Actividad sugerida}
    \begin{itemize}
      \item Capítulo 13, actividad EMO [MOV-03]: segmentar zonas por patrón de movilidad.
    \end{itemize}
  \end{block}
\end{frame}

\end{document}
